\scanpage{90}
Pour tte famille à $1$ paramètre de $D$
qui se décompose en segments
tels qu'à l'intérieur de chacun
d'eux les p.r. reste \unsure{constant} et
\deleted{\unsure{de que}} qu'aux extrémités se
trouve encore des $D_t$ qui restent
en position \unsure{standard|standart}).

On se restreindra maintenant au
cas des p.r. à une seule rive ---
i.e. telles que $\widetilde D$ soit un cercle top.
Alors l'ens. des op. de transition à
partir de $D$ est muni d'une
structure polygonale \inserted{(\unsure{c.à.d.} de deux façons, suivant convention pour
\unsure{rives} de généralisation)}. Donc le graphe \inserted{orienté}
des p.r. est polygonalisé, et \deleted{\unsure{d'où}} \unsure{donc|de même}
deux façons de définir des op. de
transport pour l'orientation. La convention
la plus intéressante semble celle où la
rive de généralisation ``tire'' sur $D$ pour génériser,
et que le transport de l'orientation sur
$\mathcal X$ soit opposé au transport parallèle de
l'orientation.

Dans ce cas, dans la carte $\mathcal X$ obtenue,
tous les sommets sont d'ordre
$4$, et ils correspondent \deleted{\unsure{ainsi}} \inserted{$1$-$1$} aux
\deleted{p.r. $D$} $(D,\deleted{\unsure{\xi}}a)$, où $D$ est
une p.r., et $a$ un ``arc d'\unsure{inertie}''
sur $\widetilde D$, de type générique i.e.
joignant deux \deleted{sommets} \inserted{pts} de $D\cap S(X_1)$
(\unsure{c.à.d.} que son \unsure{intérieur|extérieur} ne rencontre
ni d'autre élément de $S(X_1)$, ni de \unsure{rive}
de spécialisation).
La situation est surtout sympathique
si la découpe de $X$ suivant les
positions relatives envisagées est un
disque (cela peut signifier aussi
que $X$ est un plan projectif
réel, et $D$ une ps droite) ---
condition qui est invariante par
opérations élémentaires, car
alors on peut interpréter $\mathcal X$ comme
\unsure{déduite} de ces découpes par des
\unsure{recollements}, et on récupère des
\margin{89}
